\section{Conclusion}

RACK-KV addresses a three-way long-context problem: reducing historical KV storage, preserving direct access to selected history, and omitting historical contributions from attention only when that omission can be justified safely. The method responds by keeping a recent window exact, compressing older states into independently decodable anchor/residual blocks, and attaching compact geometric summaries that support query-dependent block screening.

The theoretical result is local but meaningful. It gives a rigorous upper bound on the difference between attention over the fully reconstructed cache and attention after certified skipping. Empirically, the representative-layer evaluation shows that RACK-KV can reduce serialized KV storage by roughly one third while maintaining low local attention-output error, and that no certificate violation was observed in the evaluated cases. At similar serialized-memory budgets, it is more accurate than the matched-memory selection baselines used in this study. A verified two-prompt 32-layer teacher-forced evaluation further shows that these local approximations can propagate through the full model while leaving aggregate NLL and token agreement close to the Full-KV reference on the tested scope.

What remains unresolved is equally important. Compression error is still empirical rather than rigorously certified, the current certificate is conservative, logical query-head skips do not yet imply physical GQA block omission, and the full-model evidence is still limited to two short teacher-forced prompt categories. RACK-KV should therefore be viewed as evidence that compression, random access, local certification, and narrow end-to-end quality preservation can coexist in one system, not yet as proof of broad generalization or runtime acceleration.
